\chapter{La transformada Zeta}

\section{Introducci\'{o}n}

\subsection{Origen de la transformada Zeta}

\begin{enumerate}

\item Calculemos la salida de un sistema LIT de tiempo discreto y respuesta al impulso $h[n]$ cuando la
entrada es $x[n] = z^{n}$ ($z$ es una variable compleja) mediante la suma de convoluci\'{o}n
\begin{equation}
   y[n] = \sum_{k=-\infty}^{\infty} h[k] \, x[n-k] = 
          \sum_{k=-\infty}^{\infty} h[k] \, z^{n-k}.
\end{equation}

\item Observemos que $z^{n-k} = z^{n} \, z^{-k}$. Por lo tanto podemos escribir
\begin{equation}
   y[n] = \left(\sum_{k=-\infty}^{\infty} h[k] \, z^{-k} \right) \, z^{n} = Z\{h[n]\} \, z^{n}, \label{motivozeta}
\end{equation}
dado que el factor $z^{n}$ es una constante para la sumatoria. En este caso, $Z\{h[n]\}$ representa a la
transformada $Z$ bilateral de $h[n]$.

\item En el caso especial en que el sistema sea causal, la respuesta al impulso es nula para $n<0$. 
Entonces la ecuaci\'{o}n (\ref{motivozeta}) se modifica
\begin{equation}
   y[n] = \left(\sum_{\mathbf{k=0}}^{\infty} h[k] \, z^{-k} \right) \, z^{n} = Z^{+}\{h[n]\} \, z^{n}.
\end{equation}
En este caso, $Z^{+}\{h[n]\}$ representa la transformada $Z$ unilateral de $h[n]$.
\end{enumerate}

\subsection{Ecuaci\'{o}n de an\'{a}lisis (bilateral)}
La transformada $Z$ bilateral de una secuencia $x\left[ n\right] $ se define mediante la ecuaci\'{o}n de 
an\'{a}lisis 
\begin{equation}
        X\left( z\right) =\sum_{n=-\infty }^{\infty }x\left[ n\right] \, z^{-n}.
\end{equation}

\subsection{Ecuaci\'{o}n de an\'{a}lisis (unilateral)}
La transformada $Z$ unilateral de una secuencia $x\left[ n\right] $ se define mediante la ecuaci\'{o}n de 
an\'{a}lisis 
\begin{equation}
        X^{+}\left( z\right) =\sum_{n=0}^{\infty }x\left[ n\right] \, z^{-n}.
\end{equation}
Observe que la diferencia con la transformada bilateral es que el l\'{\i}mite inferior de la sumatoria es $0$, en vez de ser $-\infty$. Esto significa que si la se\~{n}al $x[n]$ es nula para $n < 0$ ambas transformadas coinciden.

\subsection{Ecuaci\'{o}n de s\'{\i}ntesis}

La transformada $Z$ inversa (ya sea bilateral o unilateral) est\'{a} definida por la ecuaci\'{o}n de s\'{\i}ntesis 
\begin{equation}
        x\left[ n\right] =\frac{1}{2\pi j}\oint_{C}X\left( z\right) \, z^{n-1} \, dz.
\end{equation}

\subsection{Aclaraciones respecto a ambas transformadas}

\begin{enumerate}

\item La funci\'{o}n $X\left( z\right) $, tanto para el caso bilateral como el unilateral, est\'{a} 
definida para una regi\'{o}n en el plano complejo $z$, denominada \textbf{regi\'{o}n de convergencia}, 
para la cu\'{a}l la sumatoria de la ecuaci\'{o}n de an\'{a}lisis existe.

\item Como $z$ es una cantidad compleja, $X\left( z\right) $ es una funci\'{o}n compleja de una variable 
compleja. Entonces, la ecuaci\'{o}n de s\'{\i}ntesis implica integraci\'{o}n en el dominio complejo $z.$ 
Solamente usaremos esta f\'{o}rmula para probar teoremas y en algunos ejemplos puntuales.

\end{enumerate}

\subsection{Ejemplo} \label{zetaejemplouno}

\begin{enumerate}

\item Calculemos la transformada $Z$ bilateral de un impulso unitario $x[n] = \delta[n]$.
Entonces escribimos
\begin{equation}
X(z) = \sum_{k=-\infty}^{\infty} x[k] \, z^{-k} = \sum_{k=-\infty}^{\infty} \delta[k] \, z^{-k}. 
\end{equation}
Dado que $\delta[k] = 1$ si y solo si $k=0$, entonces resulta
\begin{equation}
X(z) = \sum_{k=-\infty}^{\infty} \delta[k] \, z^{-k} = \delta[0] \, z^{0} = 1. 
\end{equation}

\item Calculemos la transformada $Z$ bilateral de un impulso unitario desplazado en el tiempo $x[n] = \delta[n-m]$.
Entonces escribimos
\begin{equation}
X(z) = \sum_{k=-\infty}^{\infty} x[k] \, z^{-k} = \sum_{k=-\infty}^{\infty} \delta[k-m] \, z^{-k}. 
\end{equation}
Dado que $\delta[k-m] = 1$ si y solo si $k-m=0$, entonces resulta $k = m$ y
\begin{equation}
X(z) = \sum_{k=-\infty}^{\infty} \delta[k-m] \, z^{-k} = \delta[0] \, z^{-m} = z^{-m}. 
\end{equation}

\item Calculemos la transformada $Z$ bilateral de un escal\'{o}n unitario $x[n] = u[n]$.
Entonces escribimos
\begin{equation}
X(z) = \sum_{k=-\infty}^{\infty} x[k] \, z^{-k} = \sum_{k=-\infty}^{\infty} u[k] \, z^{-k}. 
\end{equation}
Dado que $u[k] = 1$ si y solo si $k \geq 0$, entonces resulta
\begin{equation}
X(z) = \sum_{k=-\infty}^{\infty} u[k] \, z^{-k} = \sum_{k=0}^{\infty} z^{-k}. 
\end{equation}
Observemos que $\sum_{k=0}^{\infty} z^{-k}$ es una serie geom\'{e}trica de raz\'{o}n $z^{-1}$. Para que esta serie sea convergente debe ser $|| z^{-1} || \leq 1$.

\item Calculemos la transformada $Z$ bilateral de una exponencial \textbf{causal} de duraci\'{o}n infinita
\begin{equation}
x\left[ n\right] =a^{n} \, u\left[ n\right] .
\end{equation}

\item La ecuaci\'{o}n de an\'{a}lisis correspondiente es
\begin{equation}
X\left( z\right) =\sum_{k=-\infty }^{\infty }x\left[ k\right] \, z^{-k}.
\end{equation}

\item Por lo tanto, escribimos
\begin{equation}
  X\left( z\right) = \sum_{k=-\infty}^{\infty} a^{k} \, u[k] \, z^{-k} = \sum_{k=0}^{\infty} a^k \, z^{-k} = 
    \frac{1}{1-a \, z^{-1}},
\end{equation}
donde $ \left| a \, z^{-1}\right| <1 \Leftrightarrow \left| z\right| >\left| a\right|$.

\item La regi\'{o}n de convergencia es el exterior de un c\'{\i}rculo de radio $\left| a\right| $ 
centrado en el or\'{\i}gen.

\end{enumerate}

\section{La transformada $Z$, la transformada de Fourier, y la regi\'{o}n de convergencia} \label{zetaconv}

\subsection{Desarrollo te\'{o}rico}

\begin{enumerate}
  \item Consideremos la ecuaci\'{o}n de s\'{\i}ntesis de la transformada bilateral
    \begin{equation}
       X(z) = \sum_{k=-\infty}^{\infty} x[k] \, z^{-k}.
    \end{equation}

  \item Expresemos la variable compleja $z$ en forma polar
    \begin{equation}
      z = r \, e^{j \omega},
    \end{equation}
    donde $r = \left\| z \right\|$, y $\omega = \angle(z)$,  obteniendo
    \begin{equation}
      X\left( r \, e^{j \, \omega }\right) =
        \sum_{k=-\infty }^{\infty }x\left[ k\right] \, \left( r \, e^{j \, \omega }\right) ^{-k}.
    \end{equation}
    Reagrupando factores, resulta finalmente
    \begin{equation}
      X\left( r \, e^{j \, \omega }\right) =
        \sum_{k=-\infty }^{\infty }\left( x\left[ k\right] \, r^{-k}\right) \, e^{-jk\omega }.
    \end{equation}

  \item La funci\'{o}n $X\left( r \, e^{j \, \omega }\right) $ es la Transformada de
    Fourier del producto $x\left[ k\right] \, r^{-k},$ y la Transformada Zeta
    evaluada sobre el c\'{\i}rculo unitario $e^{j \, \omega }$ coincide con la
    Transformada de Fourier. 

  \item Para que exista la transformada $Z$ de $x[k]$ debe existir la 
    transformada discreta de Fourier de $x[k] \, r^{-k}$ para alg\'{u}n valor 
    de $r$.

  \item Los valores de $r$ para los cuales existe la transformada discreta de Fourier
    de $x[k] \, r^{-k}$ definen la regi\'{o}n de convergencia de la transformada $Z$ de $x[n]$.

  \item Cuando hablamos de la existencia de $X(z)$, lo que estamos afirmando es que la magnitud de 
    $X(z)$ est\'{a} acotada por un n\'{u}mero $M$ finito
    \begin{equation}
      \left\| X(z) \right\| < M < \infty.
    \end{equation}
    Por lo tanto, para determinar la regi\'{o}n de convergencia debemos encontrar una cota superior para 
    $\left\| X(z) \right\|$. Para esto planteamos
    \begin{equation}
      \left\| X(z) \right\| = 
      \left\| \sum_{k=-\infty}^{\infty} x[k] \, r^{-k} \, e^{-j \, \omega n} \right\|.
    \end{equation}
    Aplicando la desigualdad triangular
    \begin{equation}
      \left\| X(z) \right\| \leq
      \sum_{k=-\infty}^{\infty} \left\| x[k] \, r^{-k} \, e^{-j \, \omega \, n} \right\|,
    \end{equation}
    y dado que $\left\| e^{-j \, \omega \, k} \right\| = 1$, finalmente obtenemos
    \begin{equation}
      \left\| X(z) \right\| \leq
            \sum_{k=-\infty}^{\infty} \left\| x[k] \, r^{-k} \right\| \leq M.
    \end{equation}
    
  \item Podemos sacar como conclusi\'{o}n que determinar la regi\'{o}n de convergencia equivale a hallar
    los valores de $r$ para los cuales la secuencia $x[n] \, r^{-n}$ es absolutamente sumable. Esto coincide
    con las condiciones de Dirichlet para la existencia de la transformada de Fourier.

  \item Veremos en lo sucesivo que el c\'{\i}rculo unitario en el plano $z$ equivale al eje imaginario en 
    el plano $s.$

\end{enumerate}

\subsection{Reglas pr\'{a}cticas para determinar la regi\'{o}n de convergencia}

\begin{enumerate}
\item  La regi\'{o}n de convergencia de $X\left( z\right) $ consiste en un
anillo en el plano $z$ centrado en el origen.

\item  La regi\'{o}n de convergencia no contiene polos.

\item  Si la secuencia $x\left[ n\right] $ es de duraci\'{o}n finita,
entonces la regi\'{o}n de convergencia es el plano $z$ entero, excepto
posiblemente en el origen $z=0$ y en el infinito $\left\| z \right\|=\infty$.

\item  Si $x\left[ n\right] $ es una secuencia derecha, y si el c\'{\i}rculo 
$\left\| z\right\| =r_{0}$ est\'{a} en la regi\'{o}n de convergencia,
entonces todos los valores finitos de $z$ para los cuales $\left\| z\right\|>r_{0}$ 
est\'{a}n tambi\'{e}n en la regi\'{o}n de convergencia.

\item  Si $x\left[ n\right] $ es una secuencia izquierda, y si el
c\'{\i}rculo $\left\| z\right\| =r_{0}$ est\'{a} en la regi\'{o}n de
convergencia, entonces todos los valores finitos de $z$ para los cuales 
$\left\| z\right\| <r_{0}$ est\'{a}n tambi\'{e}n en la regi\'{o}n de
convergencia.

\item  Si $x\left[ n\right] $ es una secuencia bilateral, y si el
c\'{\i}rculo $\left\| z\right\| =r_{0}$ est\'{a} en la regi\'{o}n de
convergencia, entonces la regi\'{o}n de convergencia consiste en un anillo
que contiene a $\left\| z\right\| =r_{0}$.

\item  Si $X\left( z\right) $ correspondiente a $x\left[ n\right] $ es
racional, entonces la regi\'{o}n de convergencia est\'{a} limitada por polos
o se extiende hasta el infinito.

\item  Si $X\left( z\right) $ correspondiente a $x\left[ n\right] $ es
racional, y si $x\left[ n\right] $ es derecha, entonces la regi\'{o}n de
convergencia est\'{a} ubicada fuera del c\'{\i}rculo que contiene al polo de
mayor magnitud. Adem\'{a}s, si $x\left[ n\right] $ es causal (o sea, si es
derecha y nula para $n<0),$ entonces la regi\'{o}n de convergencia
tambi\'{e}n incluye $z=\infty .$

\item  Si $X\left( z\right) $ correspondiente a $x\left[ n\right] $ es
racional, y si $x\left[ n\right] $ es izquierda, entonces la regi\'{o}n de
convergencia est\'{a} ubicada dentro del c\'{\i}rculo que contiene al polo de
menor magnitud. Adem\'{a}s, si $x\left[ n\right] $ es nocausal (o sea, si es
izquierda y nula para $n>0),$ entonces la regi\'{o}n de convergencia
tambi\'{e}n incluye $z=0$.
\end{enumerate}

\section{La transformada Zeta bilateral}

\subsection{El par funci\'{o}n temporal transformada}

Indicaremos el par formado por una funci\'{o}n temporal y su transformada como 
\begin{equation}
x\left[ n\right] 
\begin{array}{c}
Z \\ 
\Leftrightarrow
\end{array}
X\left( z\right) .
\end{equation}

\subsection{Propiedades}

La siguiente tabla resume algunas propiedades de la transformada $Z$ \textbf{bilateral}.

\begin{center}
\begin{tabular}{llll}
Se\~{n}al                  & Transformada                                  & Regi\'{o}n de convergencia \\ 
&&\\
$a \, x\left[ n\right] +b \, y\left[ n\right] $ & 
$a \, X\left( z\right) +b \, Y\left( z\right) $ & 
     $\supset \left( R_{x}\cap R_{y}\right) $ \\ 
&&\\
$x\left[ n-n_{0}\right] $               & $z^{-n_{0}} \, X\left( z\right) $           & $R_{x}^{ }$ \\ 
&&\\
$nx\left[ n\right] $                    & $-z \, \frac{dX\left(z\right) }{dz}$        & $R_{x}^{ }$ \\ 
&&\\
$x\left[ n\right] \, y\left[ n\right]$    & $X\left( z\right) * Y\left( z\right) $  & 
     $\supset \left( R_{x}\cap R_{y}\right)$                                                    \\
&&\\
$x\left[ n\right] *  y\left[ n\right]$  & $X\left( z\right) \, Y\left( z\right) $   & 
     $\supset \left( R_{x}\cap R_{y}\right)$
\end{tabular}
\end{center}

\subsection{Linealidad}

\begin{equation}
a \, x\left[ n\right] +b \, y\left[ n\right] 
\begin{array}{c}
Z \\ 
\Leftrightarrow
\end{array}
a \, X\left( z\right) +b \, Y\left( z\right) .
\end{equation}

\subsection{Demora $n_{0}$}

\begin{equation}
x\left[ n-n_{0}\right] 
\begin{array}{c}
Z \\ 
\Leftrightarrow
\end{array}
z^{-n_{0}} \, X\left( z\right) .
\end{equation}

\subsection{Multiplicaci\'{o}n por $n$}

\begin{equation}
n \, x\left[ n\right] 
\begin{array}{c}
Z \\ 
\Leftrightarrow
\end{array}
-z \, \frac{dX\left( z\right) }{dz}.
\end{equation}

\subsection{Convoluci\'{o}n en $z$}

\begin{equation}
x\left[ n\right] \,  y\left[ n\right] 
\begin{array}{c}
Z \\ 
\Leftrightarrow
\end{array}
X\left( z\right) * Y\left( z\right) .
\end{equation}

\subsection{Convoluci\'{o}n en $n$}

\begin{equation}
x\left[ n\right] * y\left[ n\right] 
\begin{array}{c}
Z \\ 
\Leftrightarrow
\end{array}
X\left( z\right) \, Y\left( z\right) .
\end{equation}

\section{Transformada Zeta de algunas secuencias}

\subsection{Secuencia impulso unitario}

\begin{equation}
x\left[ n\right] =\delta \left[ n\right] ,
\end{equation}

\begin{equation}
X\left( z\right) = \sum_{n=-\infty }^{\infty } \delta \, \left[ n\right] z^{-n},
\end{equation}

\begin{equation}
X\left( z\right) = 1 \, z^{0} = 1.
\end{equation}

Regi\'{o}n de convergencia: todo el plano $z$.

\subsection{Secuencia escal\'{o}n unitario}

\begin{equation}
x\left[ n\right] =u\left[ n\right] ,
\end{equation}

\begin{equation}
X\left( z\right) =\sum_{n=-\infty }^{\infty }u\left[ n\right] \, z^{-n},
\end{equation}

\begin{equation}
X\left( z\right) =\sum_{n=0}^{\infty }z^{-n}=\frac{1}{1-z^{-1}}, \left| z^{-1}\right| <1.
\end{equation}

Regi\'{o}n de convergencia: exterior del c\'{\i}rculo unitario centrado en el
or\'{\i}gen, $\left| z^{-1}\right| <1\Leftrightarrow \left| z\right| >1$.

\section{Relaci\'{o}n entre secuencias causales y diagramas de ceros y polos}

\subsection{Secuencia exponencial causal}

Recordemos que la transformada $Z$ de una exponencial \textbf{causal} 
\begin{equation}
x\left[ n\right] =a^{n} \, u\left[ n\right] .
\end{equation}
es
\begin{equation}
X\left( z\right) =\frac{1}{1-a \, z^{-1}},
\end{equation}
donde
\begin{equation}
 \left| a \, z^{-1}\right| <1 \Leftrightarrow \left| z\right| >\left|a\right| .
\end{equation}
y por lo tanto la regi\'{o}n de convergencia es el exterior de un c\'{\i}rculo de radio 
$\left| a\right| $ centrado en el or\'{\i}gen.

\subsection{Secuencia exponencial no causal}

\begin{enumerate}

\item Calculemos la transformada $Z$ de una exponencial \textbf{no causal}
\begin{equation}
x\left[ n\right] =-a^{n} \, u\left[ -n-1\right] .
\end{equation}

\item La ecuaci\'{o}n de an\'{a}lisis correspondiente es
\begin{equation}
X\left( z\right) =\sum_{n=-\infty }^{\infty }x\left[ n\right] \, z^{-n},
\end{equation}

\item En este caso especial, podemos escribir
\begin{equation}
X\left( z\right) = \sum_{n=-\infty }^{\infty }\left( -a^{n} \, u\left[ -n-1\right] \right) \, z^{-n}.
\end{equation}

\item Por lo tanto, resulta
\begin{equation}
X\left( z\right) =-\sum_{n=-\infty }^{-1}\left( a \, z^{-1}\right)^{n}=
  -\sum_{n=1}^{\infty }\left( a^{-1} \, z\right) ^{-n}.
\end{equation}

\item Finalmente, 
\begin{equation}
X\left( z\right) =-\sum_{n=1}^{\infty }\left( a^{-1}z\right) ^{-n}=
  \frac{\left( a^{-1} \, z\right) }{1-\left( a^{-1} \, z\right) }.
\end{equation}

\item Observe que reacomodando el resultado anterior, obtenemos
\begin{equation}
X\left( z\right) =\frac{\left( a^{-1} \, z\right) }{1-\left( a^{-1} \, z\right) }=\frac{1}{1-a \, z^{-1}},
\end{equation}
donde
\begin{equation}
 \left| a^{-1} \, z\right| <1\Leftrightarrow \left|z\right| <\left| a\right| .
\end{equation}

\item La regi\'{o}n de convergencia es el \textbf{interior} del c\'{\i}rculo de radio $\left| a\right|$.

\end{enumerate}

\subsection{Comparaci\'{o}n entre secuencias causales y no causales}

\begin{enumerate}
\item  Las secuencias causales y anticausales pueden tener la misma
transformada $Z$ pero distintas regiones de convergencia.

\item  Las secuencias causales tienen regiones de convergencia en el
exterior de un c\'{\i}rculo que intersecta al polo de la transformada $Z.$

\item  Las secuencias anticausales tienen regiones de convergencia en el
interior de un c\'{\i}rculo que intersecta al polo de la transformada $Z.$

\item  Las secuencias estables, (causales y anticausales), tienen regiones
de convergencia que incluyen el c\'{\i}rculo unitario.
\end{enumerate}

\subsection{Secuencias bilaterales}

\begin{enumerate}

\item Consideremos la secuencia bilateral 
\begin{equation}
x\left[ n\right] =a^{-n} \, u\left[ -n-1\right] +a^{n} \, u\left[ n\right] ,
\end{equation}
donde $0<a<1.$ 

\item Su transformada $Z$ es 
\begin{equation}
X\left( z\right) = 
  \begin{array}{c}
    \underbrace{\frac{1}{1-a^{-1} \, z^{-1}}} \\ 
    \left| z\right| <\frac{1}{a}
  \end{array} + 
  \begin{array}{c}
    \underbrace{\frac{1}{1-a \, z^{-1}}} \\ 
    \left| z\right| >a
  \end{array} = 
  \begin{array}{c}
    \underbrace{\frac{z \, \left( a-a^{-1}\right) }{\left( z-a^{-1}\right) \, \left(z-a\right) }} \\ 
    a<\left| z\right| <\frac{1}{a}
  \end{array}
\end{equation}

\end{enumerate}

\section{La transformada Zeta Inversa}

Veremos como determinar la secuencia a partir de su transformada Zeta.

\subsection{Primer m\'{e}todo}

\begin{enumerate}
        \item   Recordemos que la transformada Zeta de una secuencia puede ser interpretada como la 
          transformada de Fourier de una secuencia con coeficientes exponenciales
        \begin{equation}
                x[n] \, r^{-n}
                \begin{array}{l}
                        TF              \\
                        \rightarrow
                \end{array}
                X \left( r \, e^{j \omega} \right),
        \end{equation}
        para todo valor de $r$ tal que $z = r \, e^{j \omega}$ pertenezca a la regi\'{o}n de convergencia.
        
        \item   Aplicando la transformada inversa de Fourier a ambos lados de la ecuaci\'{o}n obtenemos
        \begin{equation}
                x[n] \, r^{-n} = \frac{1}{2 \pi} \int^{2 \pi}_{0} X( r \, e^{j \omega}) \, e^{j \omega n} \, d \omega,
        \end{equation}
        
        \item   Multiplicando ambos lados de la ecuaci\'{o}n por $r^{n}$ obtenemos finalmente 
        \begin{equation} \label{zetaanalisis}
                x[n] = \frac{1}{2 \, \pi} \int^{2 \pi}_{0} X( r \, e^{j \omega}) \, (r \, e^{j \omega})^{n} \, d \omega,
        \end{equation}
        Esto significa que podemos recuperar $x[n]$ partiendo de su transformada $Z$ evaluada a lo largo
        de una trayectoria $z=r \, e^{j \omega}$ inclu\'{\i}da en su regi\'{o}n de convergencia, con $r$ fija
        y $\omega$ variable en un intervalo de longitud $2 \, \pi$.

        \item Apliquemos un cambio de variable de integraci\'{o}n definido por $z = r \, e^{j \omega}$. 
          Por lo tanto
          \begin{equation}
            dz = j \, r \, e^{j \omega} \, d \omega = j \, z \, d\omega,
          \end{equation}
          y 
          \begin{equation}
            \frac{1}{j} \, z^{-1} \, dz = d\omega.
          \end{equation}
          Entonces, la integral (\ref{zetaanalisis}) se puede reescribir como
          \begin{equation}
            x[n] = \frac{1}{2 \, \pi \, j} \oint_{C} X(z) \, z^{n-1} \, dz.
          \end{equation}
          donde $C$ simboliza el contorno de integraci\'{o}n de la integral curvil\'{\i}nea.

          \item Existen dos teoremas, debidos a Cauchy, que nos permiten evaluar integrales 
          curvil\'{\i}neas en el plano complejo.

          \begin{enumerate}

             \item El primer teorema de los  residuos de Cauchy afirma que
                \begin{equation}
                   \frac{1}{2 \, \pi \, j} \, \oint_{C} \frac{f(z)}{z-z_{0}} \, dz = 
                   f(z_{0}),
                \end{equation}
                si $z_{0}$ es interior al contorno $C$, y
               \begin{equation}
                   \frac{1}{2 \, \pi \, j} \, \oint_{C} \frac{f(z)}{z-z_{0}} \, dz = 0
               \end{equation}
               si $z_{0}$ es exterior al contorno $C$, donde $f(z)$ es derivable sobre y en el interior 
               del contorno $C$ y sin polos en $z_{0}$.
         
            \item El segundo teorema de Cauchy afirma que
               \begin{equation}
                  \frac{1}{2 \, \pi \, j} \oint_{C} \frac{f(z)}{g(z)} \, dz = 
                  \sum_{i=1}^{n} A_{i}(z_{i})
               \end{equation}
               donde
               \begin{equation}
                  A_{i}(z) = (z-z_{i}) \, \frac{f(z)}{g(z)}. \label{residuozeta}
               \end{equation}
               son los residuos de $\frac{f(z)}{g(z)}$ evaluados en los polos $z_{i}$,
               si se cumplen las siguientes condiciones
               \begin{enumerate}
                  \item $f(z)$ se mantiene acotada en el interior del contorno $C$.
                  \item $g(z)$ es un polinomio con $N$ ra\'{\i}ces simples (o sea, distintas) $z_{i}$, 
                        $i \in [1,N]$  en el interior del contorno $C$.
               \end{enumerate}
          \end{enumerate}

         \item  Para poder aplicar estos teoremas, supongamos que es posible factorizar
           \begin{equation}
             X(z) \, z^{-n} = \frac{N(z)}{D(z)},
           \end{equation}
           donde $D(z)$ tiene $M$ ceros $z_{i}$ simples (distintos) en el interior del contorno $C$. 
           Entonces podemos afirmar que
           \begin{equation}
              x[n] = \sum_{i=1}^{M} A_{i}(z_{i}),
           \end{equation}
           donde (ver ec. \ref{residuozeta})
           \begin{equation}
               A_{i}(z) = (z-z_{i}) \, X(z) \, z^{-n}.
           \end{equation}

\end{enumerate}

\subsection{Segundo m\'{e}todo}

Podemos expandir $X(z)$ en una serie de potencias de $z^{-n}$
\begin{equation}
  X(z) = \sum_{n=-\infty}^{\infty} c_{n} \, z^{-n},
\end{equation}
que converge en la regi\'{o}n de convergencia correspondiente. En este caso, 
\begin{equation} x[n] = c_{n}.
\end{equation}

\subsection{Tercer m\'{e}todo}

Podemos expresar $X(z)$ como una combinaci\'{o}n lineal
\begin{equation}
  X(z) = \sum_{i=1}^{M} a_{i} \, X_{i}(z),
\end{equation}
donde conozcamos las secuencias $x_{i}[n]$ correspondientes a cada $X_{i}(z)$. En este caso
\begin{equation}
  x[n] = \sum_{i=1}^{M} a_{i} \, x_{i}[n].
\end{equation}

\subsection{Ejemplos}

\begin{enumerate}

\item Consideremos la transformada $Z$
\begin{equation}
  X(z) = 
 \frac{3 - \frac{5}{6} \, z^{-1}}{\left( 1 - \frac{1}{4} \, z^{-1} \right) \left(1 - \frac{1}{3} \, z^{-1} \right)},
\end{equation}
donde $| z | > \frac{1}{3}$.

\item Tenemos dos polos, uno en $z = \frac{1}{4}$ y otro en $z = \frac{1}{3}$. La regi\'{o}n de 
convergencia est\'{a} fuera del c\'{i}rculo definido por el polo m\'{a}s alejado del origen, y por
lo tanto (ver las reglas para determinar la regi\'{o}n de convergencia en \ref{zetaconv}) 
la secuencia es lateral derecha.

\item Expandiendo $X(z)$ en fracciones parciales tenemos
\begin{equation}
  X(z) = \frac{A}{1-\frac{1}{4} \, z^{-1}} +  \frac{B}{1-\frac{1}{3} \, z^{-1}},
\end{equation}
donde
\begin{equation}
  A = \left. X(z) \, (1-\frac{1}{4} \, z^{-1}) \right|_{z=\frac{1}{4}} = 1,
\end{equation}
y
\begin{equation}
  B = \left. X(z) \, (1-\frac{1}{3} \, z^{-1}) \right|_{z=\frac{1}{3}} = 2.
\end{equation}

\item Recordando el ejemplo desarrollado en (\ref{zetaejemplouno}) podemos deducir los siguientes pares de 
se\~{n}ales y transformadas
\begin{equation}
  x_{1}[n] = \frac{1}{4^{n}} \, u[n] 
    \begin{array}{c} Z \\ \Leftrightarrow \end{array}
  X_{1}(z) = \begin{array}{cc} \frac{1}{1-\frac{1}{4} \, z^{-1}}, & |z| > \frac{1}{4}, \end{array}
\end{equation}
\begin{equation}
  x_{2}[n] = \frac{1}{3^{n}} \, u[n] 
    \begin{array}{c} Z \\ \Leftrightarrow \end{array}
  X_{2}(z) = \begin{array}{cc} \frac{1}{1-\frac{1}{3} \, z^{-1}}, & |z| > \frac{1}{3}, \end{array}
\end{equation}
y en consecuencia, la secuencia buscada es
\begin{equation}
  x[n] = \left( \frac{1}{4^{n}} + \frac{2}{3^{n}} \right) \, u[n].
\end{equation}

\end{enumerate}

\section{Transformada Zeta Unilateral}

\subsection{Algunos comentarios}

\begin{enumerate}

\item Hasta aqu\'{\i} hemos estado trabajando con la Transformada Zeta Bilateral. Sin embargo, como en el 
caso de la Transformada de Laplace, existe una forma alternativa denominada 
\textbf{Transformada Zeta Unilateral, \'{u}til para analizar sistemas causales que no est\'{a}n en
reposo para un tiempo inicial dado}. 

\item Recordemos que la transformada $Z$ unilateral de una secuencia $x\left[ n\right] $ se define 
mediante la ecuaci\'{o}n de an\'{a}lisis 
\begin{equation}
        X^{+}\left( z\right) = Z^{+}\{x[n]\} = \sum_{n=0}^{\infty }x\left[ n\right] z^{-n}.
\end{equation}

\item La transformada unilateral $X^{+}(z)$ no contiene informaci\'{o}n de $x[n]$ para $n<0$.

\item La transformada unilateral $X^{+}(z)$ es \'{u}nica s\'{o}lo para se\~{n}ales causales.

\item La relaci\'{o}n entre las transformadas unilateral y bilateral es
\begin{equation}
  Z^{+}\left\{x[n]\right\} = Z\left\{x[n] u[n]\right\}.
\end{equation}

\item Las propiedades de la transformada $Z$ unilateral son similares a las de la transformada $Z$ 
bilateral, excepto aquellas que impliquen un desplazamiento temporal.
\end{enumerate}

\subsection{Propiedades de la transformada $Z$ unilateral}

Supondremos que para una secuencia $x[n]$ su transformada $Z$ unilateral existe
\begin{equation}
  x[n] \begin{array}{c} Z^{+} \\ \leftrightarrow \end{array} X^{+}(z).
\end{equation}

\subsubsection{Retardo de tiempo}

\begin{enumerate}

\item Si $x[n]$ no es causal, entonces
\begin{equation}
  x[n-k] 
     \begin{array}{c} Z^{+} \\ \leftrightarrow \end{array} 
  z^{-k} \, \left[ X^{+}(z) + \sum_{n=1}^{k} x[-n] \, z^{n} \right],
\end{equation}
para $k > 0$.

\item En cambio, si $x[n]$ es causal, tenemos que
\begin{equation}
  x[n-k] 
     \begin{array}{c} Z^{+} \\ \leftrightarrow \end{array} 
  z^{-k} X^{+}(z),
\end{equation}
para $k > 0$.

\end{enumerate}

\subsubsection{Avance de tiempo}

\begin{equation}
  x[n+k] 
     \begin{array}{c} Z^{+} \\ \leftrightarrow \end{array} 
  z^{k} \left[ X^{+}(z) - \sum_{n=0}^{k-1} x[n] \, z^{-n} \right],
\end{equation}
para $k > 0$.

\subsubsection{Teorema del valor final}

Si la regi\'{o}n de convergencia de $(z-1) \, X^{+}(z)$ incluye el c\'{\i}rculo unitario, entonces
\begin{equation}
  \begin{array}{c}
     Lim \\
     n \rightarrow \infty
  \end{array}
  x[n] = 
  \begin{array}{c}
     Lim \\
     z \rightarrow 1
  \end{array}
  (z - 1) \, X^{+}(z).
\end{equation}
